\documentclass[10pt, a4paper]{article}

% --- Packages ---
\usepackage[utf8]{inputenc}
\usepackage[T1]{fontenc}
\usepackage{geometry}
\geometry{top=2.5cm, bottom=2.5cm, left=2.5cm, right=2.5cm}
\usepackage{titlesec}
\usepackage{xcolor}
\usepackage{booktabs}
\usepackage{enumitem}
\usepackage{parskip}
\usepackage{hyperref}
\usepackage{mathpazo} % Palatino font for an elegant look

% --- Custom Colors ---
\definecolor{primary}{RGB}{0, 0, 0} % A professional steel blue
\definecolor{darkgray}{RGB}{70, 70, 70}

% --- Heading Formatting ---
\titleformat{\section}
  {\Large\bfseries\color{primary}}
  {\thesection}{1em}{}[\vspace{-0.5em}\rule{\textwidth}{0.4pt}]
\titlespacing*{\section}{0pt}{1.5ex plus 1ex minus .2ex}{1ex plus .2ex}

\titleformat{\subsection}
  {\large\bfseries\color{darkgray}}
  {\thesubsection}{1em}{}

% --- Document Info ---
\title{\vspace{-2cm}\textbf{\color{primary}Project Proposal: SG HomeRadar}\\
}
\author{\textbf{Group 3:} Zhao Jing, Zhao Liangyi, Liu Tianyi, Zhang Junming, Cao Tianhan}
\date{\today}

\begin{document}

\maketitle
\thispagestyle{empty}
\vspace{-0.5cm} % Reduce space after title

\section*{1. Problem Statement}
Traditional property platforms use hard filters such as price and room count, but fail to reflect neighborhood liveability, requiring users to consult maps and amenity information separately.

We aim to integrate current and historical prices, periodically refreshed listings, geospatial facilities and community data into an LLM-powered platform for efficient, personalized area and property discovery.


\section*{2. Proposed Solution}
We propose a web-based interactive recommendation system featuring two core interaction paths:
\begin{itemize}[itemsep=2pt, parsep=0pt]
    \item \textbf{Macro Exploration (Map UI):} An interactive map of Singapore's 55 URA planning areas. Hover cards show median housing prices and a "Liveability Score" based on Transport accessibility, Food and shopping, Education, Healthcare, and other factors.
    \item \textbf{Micro Retrieval (NLP Engine):} Users describe their housing needs in natural language. An LLM extracts constraints and preferences; a deterministic engine filters and ranks evidence-backed options from the knowledge base and periodically refreshed listings.
\end{itemize}

\section*{3. Data Strategy \& Web Mining}
We combine web-mined information with authoritative public datasets:
\begin{itemize}[itemsep=2pt, parsep=0pt]
    \item \textbf{Web-Mined Market Data:} Periodically refreshed property listings, infrastructure and amenities such as MRT stations, food, shopping, education and healthcare, partial sale and rental transaction records and community reviews.
    \item \textbf{Official Public Datasets:} Historical housing transactions, planning-area and subzone boundaries, addresses, and other authoritative housing and geospatial attributes.
\end{itemize}

\section*{4. Technical Architecture}
Our system consists of five main parts:

\begin{itemize}
  \item \textbf{Data and Knowledge Layer:} We regularly collect housing listings, transaction records,
  nearby facilities and community reviews. After cleaning and map location, the data is stored in a shared
  knowledge base.

  \item \textbf{Interactive Map:} Users can explore Singapore without entering a request. The map displays
  planning areas, subzones, facilities, listings, median prices, facility counts, ratings and reviews.

  \item \textbf{Recommendation Service:} Users describe their housing needs in natural language. The LLM
  identifies their budget, location, housing type and other preferences. The system applies hard constraints
  from their requirements, such as a maximum price, then ranks the remaining options using LLM-identified
  preferences and a predefined weighted scoring function. Finally, a structured recommendation list is generated, with each recommendation explained by LLM based on the knowledge base.

  \item \textbf{Data Updates and System Connection:} Both functions use the same knowledge base.
  Housing listings are updated periodically, while other data is cleaned and processed regularly.

  \item \textbf{Evaluation:} We will test the stability and accuracy of the filtering and ranking systems,
  compare predicted prices with observed listing prices, and conduct a small user study.
\end{itemize}

\end{document}
